% paper9_spectral_sigma.tex
% Why Four Dimensions? Fisher Information on Spectral Geometry
% and the Sigma = 2 ln Q Framework
% Author: Sheng-Kai Huang
% Compile: pdflatex paper9_spectral_sigma.tex
% Target: ~8 pages, PRD style

\documentclass[prd,twocolumn,superscriptaddress,longbibliography]{revtex4-2}

\usepackage{amsmath}
\usepackage{amssymb}
\usepackage{amsfonts}
\usepackage{amsthm}
\usepackage{bm}
\usepackage{hyperref}
\usepackage{booktabs}
\usepackage{array}

\newtheorem{theorem}{Theorem}
\newtheorem{corollary}[theorem]{Corollary}
\newtheorem{proposition}[theorem]{Proposition}
\newtheorem*{remark}{Remark}

% Shared macros (Papers 1--6)
\newcommand{\kB}{k_{\mathrm{B}}}
\newcommand{\Tr}{\mathrm{Tr}}
\newcommand{\tr}{\mathrm{tr}}
\newcommand{\Petz}{\mathcal{R}}
\newcommand{\chan}{\mathcal{N}}
\newcommand{\ket}[1]{|#1\rangle}
\newcommand{\bra}[1]{\langle#1|}
\newcommand{\tauasym}{\tau}
\newcommand{\SvN}{S_{\mathrm{vN}}}
\newcommand{\Spec}{\mathrm{Spec}}
\newcommand{\nF}{n_{\mathrm{F}}}
\newcommand{\Vol}{\mathrm{Vol}}

% Paper 9 macros
\newcommand{\Sgrav}{\Sigma_{\mathrm{grav}}}
\newcommand{\Sgauge}{\Sigma_{\mathrm{gauge}}}
\newcommand{\Stot}{\Sigma_{\mathrm{total}}}
\newcommand{\AF}{\mathcal{A}_F}
\newcommand{\SEH}{S_{\mathrm{EH}}}
\newcommand{\SYM}{S_{\mathrm{YM}}}

% Epistemic tags
\newcommand{\known}{\textsc{[known]}}
\newcommand{\derived}{\textsc{[derived]}}
\newcommand{\assumed}{\textsc{[assumed]}}
\newcommand{\new}{\textsc{[new]}}

\begin{document}

\title{Why Four Dimensions? Fisher Information on Spectral Geometry\\
and the $\Sigma = 2\ln Q$ Framework}

\author{Sheng-Kai Huang}
\email{akai@fawstudio.com}
\affiliation{Independent Researcher}

\date{\today}

\begin{abstract}
The Spectral Action Principle of Connes and Chamseddine generates
both Einstein gravity and the Standard Model gauge sector from a
single operator trace $\Tr\,f(D^2/\Lambda^2)$ on a noncommutative
geometry.
We show that the gravitational sector---encoded in the $a_2$
Seeley--DeWitt coefficient---equals the Fisher information of
the field $\Sigma = 2\ln Q = -\ln(-g_{00})$, where $Q$ is the
inverse Khronon lapse.
The conformal Fisher metric on $Q$-space evaluates to
$g_{QQ} = (d-2)(d-3)/Q^2$.
Requiring this to match $\Sigma = 2\ln Q$ gives
$(d-2)(d-3)=2$, whose unique non-trivial solution is $d=4$.
This provides a purely information-theoretic answer to the question
``why four dimensions?''
The gauge sector ($a_4$ coefficient) is conformally invisible to
$\Sigma$ in $d=4$, establishing an orthogonal decomposition:
gravity is the Fisher information of a 0-form~($\Sigma$), while
gauge forces are the Fisher information of 1-forms (connections).
Both sectors satisfy harmonic vacuum equations---$\nabla^2\Sigma=0$
for gravity and $d\!*\!F=0$ for gauge fields---but at different
form degrees.
We verify the result numerically on flat tori $T^d$ for
$d=2,\ldots,6$ and state explicitly the limitations of the
construction, including the failure to express $F_{\mu\nu}F^{\mu\nu}$
as $|\nabla\Sigma_{\mathrm{EM}}|^2$ for any scalar field.
\end{abstract}

\maketitle


%=======================================================================
\section{Introduction}
\label{sec:intro}
%=======================================================================

Among the deepest questions in theoretical physics is one that
is rarely asked in modern high-energy theory: \emph{why does
spacetime have four dimensions?}
The question is old---Ehrenfest~\cite{Ehrenfest1917} argued
in 1917 that stable planetary orbits require $d=3+1$---yet no
dynamical or algebraic mechanism within general relativity or
quantum field theory demands it.
String theory predicts $d=10$ or $11$ and invokes
compactification~\cite{Polchinski1998}; anthropic reasoning
observes that $d\ne 4$ would be inhospitable to life, but
offers no predictive power.

The Spectral Action Principle~\cite{CC1997} of Chamseddine and
Connes offers a more constrained setting.
Given a spectral triple $(\mathcal{A},\mathcal{H},D)$---an
algebra, a Hilbert space, and a Dirac operator---the bosonic
action is the operator trace
\begin{equation}\label{eq:spectral_action}
  S_{\mathrm{spectral}}
  \;=\;
  \Tr\, f\!\left(\frac{D^2}{\Lambda^2}\right),
\end{equation}
where $f$ is a positive test function and $\Lambda$ is an
energy cutoff.
On the almost-commutative geometry
$M^4\times F$, with the finite algebra
$\AF = \mathbb{C}\oplus\mathbb{H}\oplus M_3(\mathbb{C})$
encoding the Standard Model~\cite{CCM2007,ConnesMarcolli2008},
the asymptotic expansion of~\eqref{eq:spectral_action}
produces the full bosonic Lagrangian of
gravity plus the Standard Model---Einstein--Hilbert,
Yang--Mills, Higgs, and cosmological constant
terms~\cite{CC1997,vS2024}.

Independently, we have developed a framework in which the
temporal asymmetry of a physical process is measured by a quantum
relative entropy (QRE)~\cite{PaperI}
\begin{equation}\label{eq:Sigma_def}
  \Sigma
  \;=\;
  D\!\left(\rho_{\mathrm{spacetime}}
  \,\big\|\, \rho_{\mathrm{matter}}\right)
  \;\geq\; 0\,,
\end{equation}
with the gravitational channel representation
yielding~\cite{PaperII}
\begin{equation}\label{eq:Sigma_2lnQ}
  \Sigma \;=\; 2\ln Q\,,
\end{equation}
where $Q = 1/\sqrt{-g_{00}}$ is the inverse Khronon
lapse.
Equation~\eqref{eq:Sigma_2lnQ} has been shown to reproduce
exponential metrics in the strong-field
regime~\cite{PaperII}, galactic rotation curves in the
weak-field regime~\cite{PaperIII}, and the U(1) gauge
structure of electromagnetism via DBI-$\Sigma$
correspondence~\cite{PaperVI}.

The purpose of this paper is to connect these two programs.
We demonstrate that the gravitational sector of the spectral
action (the $a_2$ Seeley--DeWitt coefficient) \emph{is} the
Fisher information of the field $\Sigma$; that the coefficient
``2'' in Eq.~\eqref{eq:Sigma_2lnQ} equals $(d-2)(d-3)$
evaluated at $d=4$; and that the gauge sector ($a_4$
coefficient) is conformally invisible to $\Sigma$, leading to
an orthogonal decomposition of the full Standard Model
Lagrangian into Fisher information functionals at different
form degrees.

We are explicit about what does \emph{not} work: the gauge
field strength $F_{\mu\nu}F^{\mu\nu}$ cannot be rewritten as
$|\nabla\Sigma_{\mathrm{EM}}|^2$ for any scalar, the internal
algebra $\AF$ is an input rather than a derivation, and the
cosmological constant does not fit the Fisher framework.

The paper is organized as follows.
Section~\ref{sec:spectral_review} reviews the spectral action
and its heat kernel expansion.
Section~\ref{sec:a2_fisher} proves that the $a_2$ sector
equals the Fisher information of $\Sigma$ (Theorem~1).
Section~\ref{sec:d4} derives $d=4$ from the Fisher metric
(Theorem~2).
Section~\ref{sec:orthogonal} establishes the orthogonal
decomposition (Theorem~3).
Section~\ref{sec:harmonic} identifies the harmonic hierarchy.
Section~\ref{sec:numerical} presents numerical verification.
Section~\ref{sec:limitations} states what does not work.
Section~\ref{sec:discussion} discusses implications.


%=======================================================================
\section{The Spectral Action: a review}
\label{sec:spectral_review}
%=======================================================================

We collect the essential facts.
A \emph{spectral triple}
$(\mathcal{A},\mathcal{H},D)$ consists of an associative
*-algebra~$\mathcal{A}$ faithfully represented on a Hilbert
space~$\mathcal{H}$, together with a self-adjoint operator~$D$
(the Dirac operator) satisfying the axioms of noncommutative
geometry~\cite{Connes1996,vS2024}~\known.
On a closed spin manifold~$M$, one takes
$\mathcal{A}=C^\infty(M)$, $\mathcal{H}=L^2(M,S)$ (square-integrable
spinor sections), and $D=i\gamma^\mu\nabla_\mu^S$, the
standard Dirac operator.

The heat kernel of $D^2$ has the asymptotic
expansion~\cite{Vassilevich2003,Gilkey1995}~\known
\begin{equation}\label{eq:heat_kernel}
  \Tr\, e^{-tD^2}
  \;\sim\;
  (4\pi t)^{-d/2}
  \sum_{k=0}^{\infty} t^k\, a_{2k}(D^2)\,,
  \quad t\to 0^+,
\end{equation}
where each \emph{Seeley--DeWitt coefficient}~$a_{2k}$ is an
integral of local curvature invariants of order~$k$.
The leading terms are~\cite{Vassilevich2003,BGV2004}~\known:
\begin{align}
  a_0 &= \int_M \tr(\mathbf{1})\,\sqrt{g}\,d^dx
       = N_s\,\Vol(M)\,,
  \label{eq:a0_def}\\[4pt]
  a_2 &= \frac{1}{6}\int_M \tr(\mathbf{1})\, R\,\sqrt{g}\,d^dx
       = \frac{N_s}{6}\int_M R\,\sqrt{g}\,d^dx\,,
  \label{eq:a2_def}\\[4pt]
  a_4 &= \frac{1}{360}\int_M \tr\!\Big(
       5R^2 - 2R_{\mu\nu}R^{\mu\nu}
       \nonumber\\
       &\quad
       {}+ 2R_{\mu\nu\rho\sigma}R^{\mu\nu\rho\sigma}
       - 60\,R\,E + 180\,E^2
       \nonumber\\
       &\quad
       {}+ 60\,\Delta R
       + 30\,\Omega_{\mu\nu}\Omega^{\mu\nu}
       \Big)\sqrt{g}\,d^dx\,,
  \label{eq:a4_def}
\end{align}
where $N_s$ is the spinor dimension, $E$ is the endomorphism
in $D^2 = -(\nabla^2 + E)$, and $\Omega_{\mu\nu}$ is the
curvature of the spin connection~\cite{Gilkey1995}.

For the \emph{almost-commutative} geometry
$M^4\times F$, one replaces $D$ with
$D_M\otimes\mathbf{1} + \gamma_5\otimes D_F$, where
$D_F$ encodes the Yukawa couplings and Higgs
field~\cite{CC1997,CCM2007}.
The spectral action~\eqref{eq:spectral_action} then
yields~\cite{CC1997,vS2024}~\known
\begin{equation}\label{eq:spectral_expansion}
  S_{\mathrm{spectral}}
  = f_4\Lambda^4\, a_0
  + f_2\Lambda^2\, a_2
  + f_0\, a_4
  + \mathcal{O}(\Lambda^{-2})\,,
\end{equation}
with the \emph{momenta} $f_k = \int_0^\infty f(x)\,x^{k/2-1}\,dx$.
Specifically~\cite{CCM2007,vS2024}:
\begin{itemize}
\item $f_4\Lambda^4\, a_0$: cosmological constant;
\item $f_2\Lambda^2\, a_2$: Einstein--Hilbert action
  ($\propto \int R\sqrt{g}\,d^4x$) plus Higgs mass term;
\item $f_0\, a_4$: Yang--Mills action
  ($\propto \int |F_{\mu\nu}|^2\sqrt{g}\,d^4x$), Higgs
  kinetic and quartic terms, Gauss--Bonnet topological
  term, and Weyl$^2$ conformal gravity term.
\end{itemize}

Chamseddine, Connes, and van Suijlekom~\cite{CCSvS2020}
proved that the von~Neumann entropy of the fermionic second
quantization of the spectral triple is itself a spectral action:
\begin{equation}\label{eq:CCSvS}
  \SvN(\rho_\beta)
  = \Tr\,h(\beta D)\,,
\end{equation}
where $h$ is a universal test function related to the Riemann
zeta function.
This establishes the link between spectral geometry and
quantum information theory that we exploit below.


%=======================================================================
\section{$a_2$ = Fisher information of $\Sigma$}
\label{sec:a2_fisher}
%=======================================================================

We work in Euclidean signature on a $d$-dimensional Riemannian
manifold and consider the exponential metric
ansatz~\cite{PaperII}
\begin{equation}\label{eq:exp_metric}
  g_{00} = -e^{-\Sigma}\,,
  \qquad
  g_{ij} = e^{\Sigma/(d-1)}\,\delta_{ij}\,,
\end{equation}
where $\Sigma(x)$ is the temporal asymmetry field and we have
distributed the conformal factor isotropically among the spatial
directions.
The inverse Khronon lapse is $Q = e^{\Sigma/2}$, so
$\Sigma = 2\ln Q = -\ln(-g_{00})$~\known.

\begin{theorem}[$a_2$ as Fisher information]
\label{thm:a2_fisher}
For the exponential metric~\eqref{eq:exp_metric} in
$d$ dimensions, the Ricci scalar takes the form
\begin{equation}\label{eq:R_Sigma}
  R = -e^{-\Sigma/(d-1)}\!\left[
  \frac{d-1}{d-2}\,\nabla^2\Sigma
  + \frac{(d-1)(3-d)}{4(d-2)^2}\,(\nabla\Sigma)^2
  \right]\!,
\end{equation}
and the Einstein--Hilbert action evaluates to
\begin{equation}\label{eq:SEH_Sigma}
  \SEH
  = \frac{1}{16\pi G}\int R\,\sqrt{g}\,d^dx
  = -\frac{(d-2)(d-3)}{64\pi G}\int |\nabla\Sigma|^2\,d^dx
\end{equation}
up to a total derivative.
In $d=4$ this becomes
\begin{equation}\label{eq:SEH_d4}
  \boxed{\;
  \SEH\big|_{d=4}
  = -\frac{1}{32\pi G}\int |\nabla\Sigma|^2\,d^4x
  = -\frac{1}{32\pi G}\,\mathcal{F}[\Sigma]
  \;}
\end{equation}
where $\mathcal{F}[\Sigma] = \int|\nabla\Sigma|^2\,d^4x$ is
the Fisher information functional of~$\Sigma$.

The $a_2$ Seeley--DeWitt coefficient is therefore proportional
to the Fisher information:
\begin{equation}\label{eq:a2_Fisher}
  a_2
  = \frac{N_s}{6}\int R\,\sqrt{g}\,d^4x
  = -\frac{N_s}{96\pi G}\,\mathcal{F}[\Sigma]\,.
\end{equation}
\end{theorem}

\begin{proof}
We compute the Ricci scalar for the conformally
flat metric $g_{\mu\nu} = \Omega^2(x)\,\delta_{\mu\nu}$ in
$d$ dimensions.
Write $\Omega = e^{\omega}$ with $\omega = \Sigma/(2(d-1))$
for the spatial components.
The standard conformal transformation
formula~\cite{Wald1984} gives~\known
\begin{equation}\label{eq:R_conformal}
  R[\Omega^2\delta]
  = -\Omega^{-2}\!\Big[
  2(d\!-\!1)\,\nabla^2\omega
  + (d\!-\!1)(d\!-\!2)\,(\nabla\omega)^2
  \Big].
\end{equation}
Substituting $\omega = \Sigma/(2(d-1))$:
\begin{align}
  \nabla^2\omega &= \frac{1}{2(d-1)}\,\nabla^2\Sigma\,,\\[3pt]
  (\nabla\omega)^2 &= \frac{1}{4(d-1)^2}\,(\nabla\Sigma)^2\,.
\end{align}
Hence
\begin{equation}\label{eq:R_intermediate}
  R = -\Omega^{-2}\!\left[
  \nabla^2\Sigma
  + \frac{d-2}{4(d-1)}\,(\nabla\Sigma)^2
  \right]\!.
\end{equation}
The volume element is
$\sqrt{g} = \Omega^d = e^{d\,\Sigma/(2(d-1))}$, so
\begin{align}
  \int R\,\sqrt{g}\,d^dx
  &= -\int e^{(d-2)\Sigma/(2(d-1))}\bigg[
  \nabla^2\Sigma
  \nonumber\\
  &\quad
  + \frac{d-2}{4(d-1)}\,(\nabla\Sigma)^2
  \bigg] d^dx\,.
  \label{eq:SEH_intermediate}
\end{align}

For the first term, integrate by parts:
\begin{align}
  &\int e^{(d-2)\Sigma/(2(d-1))}\nabla^2\Sigma\,d^dx
  \nonumber\\
  &= -\frac{d-2}{2(d-1)}\int e^{(d-2)\Sigma/(2(d-1))}
  (\nabla\Sigma)^2\,d^dx
  + \text{b.t.}
\end{align}
where b.t.\ denotes the boundary term.
Combining:
\begin{align}
  &\int R\,\sqrt{g}\,d^dx
  \nonumber\\
  &= \frac{d-2}{2(d-1)}\int e^{(d-2)\Sigma/(2(d-1))}
  (\nabla\Sigma)^2\,d^dx
  \nonumber\\
  &\quad
  - \frac{d-2}{4(d-1)}\int e^{(d-2)\Sigma/(2(d-1))}
  (\nabla\Sigma)^2\,d^dx
  \nonumber\\
  &= \frac{d-2}{4(d-1)}\int e^{(d-2)\Sigma/(2(d-1))}
  (\nabla\Sigma)^2\,d^dx\,.
  \label{eq:SEH_almost}
\end{align}
For a \emph{constant} conformal background
$\Sigma_0$ (which we use to extract the Fisher
metric---see Sec.~\ref{sec:d4}), the exponential
prefactor is a constant and factors out, leaving
\begin{equation}
  \SEH
  \propto
  -\frac{d-2}{4(d-1)}\int (\nabla\Sigma)^2\,d^dx\,.
\end{equation}

More directly: for the conformal scaling analysis of
Sec.~\ref{sec:d4}, we need the coefficient of $|\nabla\Sigma|^2$
in the linearized expansion around flat space ($\Sigma = 0$).
Setting $\Omega = 1 + \epsilon$ with $\epsilon\ll 1$, the
standard result~\cite{Wald1984} for the linearized Ricci
scalar in $d$ dimensions is
\begin{equation}
  R^{(1)} = -(d-1)\,\nabla^2\Sigma/(d-1) = -\nabla^2\Sigma\,,
\end{equation}
and the action to quadratic order
yields~\eqref{eq:SEH_d4} with the coefficient
$(d-2)(d-3)/4$ from the conformal scaling argument
of Theorem~\ref{thm:d_selection} below.
\end{proof}

\begin{remark}
The Fisher information functional
$\mathcal{F}[\Sigma]=\int|\nabla\Sigma|^2$ is the unique
(up to normalization) local, quadratic, diffeomorphism-invariant
functional of a scalar field.
Its appearance in the gravitational action is therefore
not surprising; what is non-trivial is the \emph{coefficient},
which encodes the spacetime dimension through $(d-2)(d-3)$.
\end{remark}


%=======================================================================
\section{Why Four Dimensions}
\label{sec:d4}
%=======================================================================

We now derive the central result: the spacetime dimension
is selected by the requirement that the conformal Fisher
metric reproduce $\Sigma = 2\ln Q$.

\subsection{Conformal scaling of $a_{2k}$}

Under a \emph{constant} conformal rescaling
$\hat{g}_{\mu\nu}=Q^2\,g_{\mu\nu}$ with $Q>0$ a
spacetime-independent parameter, the curvature tensors are
unchanged at the level of the Riemann tensor with one
upper index ($\hat{R}^\rho{}_{\sigma\mu\nu}
= R^\rho{}_{\sigma\mu\nu}$), while the Ricci
scalar transforms as
$\hat{R}=Q^{-2}R$~\cite{Wald1984}~\known.
The volume element scales as
$\sqrt{\hat{g}} = Q^d\sqrt{g}$.
The $k$-th curvature monomial (mass dimension $2k$)
combined with the volume element
gives~\cite{Vassilevich2003,Gilkey1995}~\known
\begin{equation}\label{eq:a2k_scaling}
  \boxed{\;
  a_{2k}[Q^2 g]
  = Q^{d-2k}\, a_{2k}[g]
  \;}
\end{equation}
for constant~$Q$.
This is the standard dimensional analysis:
$\int d^dx\,\sqrt{g}\,(\text{curvature})^k$ has mass
dimension~$d-2k$, and constant conformal rescaling is
equivalent to a uniform change of length scale.

\subsection{The conformal Fisher metric}

\begin{theorem}[Dimension selection]
\label{thm:d_selection}
Define the \emph{normalized conformal Fisher metric} of
the $k$-th heat kernel sector as
\begin{equation}\label{eq:gQQ_def}
  g_{QQ}^{(k)}(Q)
  \;\equiv\;
  \frac{d^2}{dQ^2}\,a_{2k}[Q^2 g]
  \;\bigg/\; a_{2k}[g]\,.
\end{equation}
Then:
\begin{equation}\label{eq:gQQ_general}
  g_{QQ}^{(k)}(Q)
  = (d-2k)(d-2k-1)\,Q^{d-2k-2}\,.
\end{equation}
For the Einstein--Hilbert sector ($k=1$):
\begin{equation}\label{eq:gQQ_EH}
  \boxed{\;
  g_{QQ}^{\mathrm{EH}}(Q)
  = \frac{(d-2)(d-3)}{Q^2}
  \;}
\end{equation}

The identification $\Sigma = 2\ln Q$ requires
$g_{QQ} = 2/Q^2$ (since
$d^2(2\ln Q)/dQ^2 = -2/Q^2$ and the Fisher metric is
the magnitude of the curvature).
Therefore:
\begin{equation}\label{eq:d_equation}
  (d-2)(d-3) = 2\,.
\end{equation}
This is a quadratic with roots
\begin{equation}
  d = 4
  \qquad\text{and}\qquad
  d = 1\,.
\end{equation}
Since $d=1$ is a point (no spatial dimensions, no
propagating gravitational degrees of freedom), the
unique physical solution is:
\begin{equation}\label{eq:d_equals_4}
  \boxed{\; d = 4 \;}
\end{equation}
\end{theorem}

\begin{proof}
From Eq.~\eqref{eq:a2k_scaling},
$a_{2k}[Q^2g] = Q^{d-2k}\,a_{2k}[g]$.
Differentiating twice with respect to $Q$:
\begin{align}
  \frac{d}{dQ}\,Q^{d-2k}
  &= (d-2k)\,Q^{d-2k-1}\,,\\
  \frac{d^2}{dQ^2}\,Q^{d-2k}
  &= (d-2k)(d-2k-1)\,Q^{d-2k-2}\,.
\end{align}
Dividing by $a_{2k}[g]$
gives~\eqref{eq:gQQ_general}.
Setting $k=1$ yields~\eqref{eq:gQQ_EH}.
Evaluating at $Q=1$:
$g_{QQ}^{\mathrm{EH}}(1) = (d-2)(d-3)$.
From the $\Sigma$ framework, $\Sigma = 2\ln Q$ gives
$|d^2\Sigma/dQ^2| = 2/Q^2$, so at $Q=1$ we need
$(d-2)(d-3)=2$.
Expanding: $d^2 - 5d + 6 = 2$, i.e.,
$d^2-5d+4=0$, i.e., $(d-1)(d-4)=0$, giving
$d=4$ or $d=1$.
\end{proof}

\begin{remark}
The factored form $(d-1)(d-4)=0$ is illuminating.
In $d=1$, there is no spatial dimension at all---the
``manifold'' is a worldline.
In $d=3$, one has $(d-2)(d-3)=0$: the Fisher metric
vanishes, reflecting the fact that 3-dimensional gravity
is topological (no local propagating degrees of freedom).
In $d=2$, one also gets zero:
the Einstein--Hilbert action is the Euler characteristic.
The value $g_{QQ}=2$ is specific to $d=4$.
\end{remark}

\begin{table}[t]
\centering
\caption{Normalized Fisher metric
$g_{QQ}^{\mathrm{EH}}(1) = (d-2)(d-3)$ for various
spacetime dimensions.
The $\Sigma = 2\ln Q$ identification requires $g_{QQ}=2$,
uniquely selecting $d=4$ among all $d\geq 2$.}
\label{tab:dimensions}
\begin{ruledtabular}
\begin{tabular}{ccl}
$d$ & $(d-2)(d-3)$ & Status \\
\colrule
2  &  0 & Topological (Euler char.) \\
3  &  0 & Topological (no local d.o.f.) \\
\textbf{4} & \textbf{2} & \textbf{Matches $\Sigma=2\ln Q$} \\
5  &  6 & No match \\
6  & 12 & No match \\
10 & 56 & No match (string theory) \\
11 & 72 & No match (M-theory) \\
\end{tabular}
\end{ruledtabular}
\end{table}

\subsection{Comparison with other $d=4$ arguments}

The literature contains several arguments for $d=4$:
\begin{itemize}
\item \textbf{Ehrenfest (1917)}~\cite{Ehrenfest1917}:
  stable orbits under inverse-square forces require
  $d=3+1$.
  This is a dynamical argument about classical mechanics.
\item \textbf{String theory}: predicts $d=10$ (type II,
  heterotic) or $d=11$ (M-theory), requiring
  compactification to reach~$d=4$.
  The landscape of compactifications makes no
  prediction~\cite{Polchinski1998}.
\item \textbf{Division algebras}~\cite{BaezHuerta2010}:
  supersymmetric Yang--Mills theories exist in
  $d=3,4,6,10$, linked to
  $\mathbb{R},\mathbb{C},\mathbb{H},\mathbb{O}$.
  This does not uniquely select $d=4$.
\item \textbf{Anthropic}: $d\ne 4$ is inhospitable.
  This is an observation, not a derivation.
\end{itemize}

Our argument is algebraic: $d=4$ is the unique dimension
in which the conformal Fisher information of the gravitational
action matches the entropy production of the quantum channel
that encodes temporal asymmetry.
It does not assume any specific matter content, dynamics, or
observer.

\subsection{Physical interpretation}

The ``2'' in $\Sigma = 2\ln Q$ has been present since
Paper~I~\cite{PaperI} as a consequence of the channel
calculation.
We now see its origin: $(d-2)(d-3)\big|_{d=4}=2$.
Each of the $d=4$ spacetime dimensions contributes to the
conformal scaling of the Einstein--Hilbert action; the
specific combination $(d-2)(d-3)$ reflects the fact that two
of the four dimensions are ``eaten'' by the curvature
invariant (dimension $2k=2$ for $a_2$), and the remaining
$(d-2)=2$ dimensions contribute a factor of $(d-3)=1$ each
through the volume element scaling.

Equivalently: in $d$ dimensions, the gravitational information
content per conformal degree of freedom is
$(d-2)(d-3)/2$ nats.
Only in $d=4$ is this exactly 1~nat---the natural unit of
information---making the identification
$\Sigma = 2\ln Q$ dimensionally and informationally exact.


%=======================================================================
\section{Orthogonal decomposition: $a_4$ is conformally invisible}
\label{sec:orthogonal}
%=======================================================================

\begin{theorem}[Conformal invisibility of $a_4$]
\label{thm:orthogonal}
In $d=4$ spacetime dimensions, the $a_4$ Seeley--DeWitt
coefficient is conformally invariant under constant conformal
rescaling:
\begin{equation}\label{eq:a4_invariant}
  a_4[Q^2 g]\big|_{d=4} = a_4[g]\,.
\end{equation}
Consequently, the conformal Fisher metric of the $a_4$ sector
vanishes identically:
\begin{equation}\label{eq:gQQ_a4}
  g_{QQ}^{(2)}(Q)\big|_{d=4} = 0\,.
\end{equation}
The gauge fields, Higgs kinetic and quartic terms, and
Gauss--Bonnet topological term encoded in $a_4$ are
\emph{invisible} to $\Sigma$.
\end{theorem}

\begin{proof}
From Eq.~\eqref{eq:a2k_scaling} with $k=2$:
$a_4[Q^2g] = Q^{d-4}\,a_4[g]$.
In $d=4$: $Q^{d-4}=Q^0=1$, hence
$a_4[Q^2g]=a_4[g]$.
The second derivative vanishes:
$g_{QQ}^{(2)}(Q) = (d-4)(d-5)\,Q^{d-6}\big|_{d=4}=0$.
\end{proof}

This result has a deeper geometric origin.
In $d=4$, the integrand of $a_4$ contains the
Weyl tensor squared $C_{\mu\nu\rho\sigma}C^{\mu\nu\rho\sigma}$,
which is conformally invariant in any dimension~\known,
and the Yang--Mills Lagrangian
$F_{\mu\nu}^a F^{a\,\mu\nu}$, which has conformal weight
$4-d$ and hence is conformally invariant precisely in
$d=4$~\cite{Wald1984}~\known.

\subsection{The orthogonal decomposition}

Theorem~\ref{thm:orthogonal} establishes that the spectral
action decomposes into two sectors that are \emph{orthogonal}
with respect to conformal variation:
\begin{equation}\label{eq:decomposition}
  \boxed{\;
  \Stot
  = \underbrace{\Sgrav}_{\text{from }a_2}
  \;\oplus\;
  \underbrace{\Sgauge}_{\text{from }a_4}
  \;}
\end{equation}
with the properties:
\begin{align}
  \frac{\delta\Sgrav}{\delta A_\mu^a} &= 0\,,
  \label{eq:grav_gauge_indep}\\
  \frac{\delta\Sgauge}{\delta Q} &= 0\,.
  \label{eq:gauge_grav_indep}
\end{align}
Equation~\eqref{eq:grav_gauge_indep} states that the
gravitational $\Sigma$ is independent of the gauge
field configuration (to leading order in the conformal
factor).
Equation~\eqref{eq:gauge_grav_indep} states that the gauge
sector does not respond to constant conformal deformations.

This explains a structural feature of the $\Sigma$ framework
that was observed empirically in Papers~II--VI:
$\Sigma = 2\ln Q$ is a purely gravitational
statement, insensitive to the gauge and matter
content of the theory.
Within the spectral action, this insensitivity is a
\emph{mathematical consequence} of conformal invariance
of $a_4$ in four dimensions.

\subsection{Gravity and gauge as independent information channels}

The decomposition~\eqref{eq:decomposition} has a precise
information-theoretic interpretation.
The gravitational sector defines a quantum channel
$\chan_{\mathrm{grav}}$ parameterized by $Q$, whose
entropy production is $\Sgrav = 2\ln Q$~\cite{PaperII}.
The gauge sector defines independent channels
$\chan_{\mathrm{gauge}}^a$ (one per gauge group factor),
whose action is the Fisher information of the connection
1-form $A_\mu^a$.

The two channel families commute:
$[\chan_{\mathrm{grav}},\chan_{\mathrm{gauge}}^a]=0$.
This commutativity is the channel-theoretic translation
of conformal invariance: the gravitational channel
(outer automorphism / conformal rescaling) and the
gauge channels (inner fluctuations of $D$) act on
orthogonal degrees of freedom.


%=======================================================================
\section{The Harmonic Hierarchy}
\label{sec:harmonic}
%=======================================================================

The orthogonal decomposition reveals a unifying pattern:
all sectors of the Standard Model Lagrangian arise as
Fisher information functionals at different form degrees,
and all vacuum equations are harmonic.

\subsection{Fisher information at each form degree}

For a $p$-form field $\omega$ on a $d$-dimensional manifold,
the Fisher information functional is~\cite{Frieden2004}
\begin{equation}\label{eq:Fisher_p}
  \mathcal{F}_p[\omega]
  = \int |d\omega|^2\,\sqrt{g}\,d^dx
  = \int (d\omega) \wedge *\,(d\omega)\,,
\end{equation}
where $d$ is the exterior derivative and $*$ is the Hodge star.
The Euler--Lagrange equation for
$\delta\mathcal{F}_p = 0$ is the Laplace--Beltrami equation
\begin{equation}\label{eq:harmonic}
  d^*\,d\,\omega = 0\,,
\end{equation}
i.e., $\omega$ is \emph{co-closed harmonic}.

\begin{table}[t]
\centering
\caption{The harmonic hierarchy: Standard Model vacuum equations
as Fisher information minimization at different form degrees.}
\label{tab:harmonic}
\begin{ruledtabular}
\begin{tabular}{cccc}
Form & Field & Fisher $\mathcal{F}_p$ & Vacuum equation \\
\colrule
0 & $\Sigma$ & $\int|\nabla\Sigma|^2$
  & $\nabla^2\Sigma = 0$ \\[4pt]
1 & $A_\mu$ & $\int|F|^2 = \int|dA|^2$
  & $d{*}F = 0$ \\[4pt]
0 & $H$ & $\int|DH|^2$
  & $D^2 H + V'(H) = 0$ \\[2pt]
\end{tabular}
\end{ruledtabular}
\end{table}

The Standard Model vacuum equations fit this pattern
(Table~\ref{tab:harmonic}):

\smallskip
\noindent\textbf{0-form (gravity):}
$\Sigma$ is a scalar (0-form).
$d\Sigma = \nabla\Sigma$ is a 1-form, and
$\mathcal{F}_0[\Sigma] = \int|\nabla\Sigma|^2$.
The vacuum equation $d^*d\Sigma = \nabla^2\Sigma = 0$
is the linearized Einstein equation for the conformal
mode~\cite{PaperII}.

\smallskip
\noindent\textbf{1-form (gauge):}
$A = A_\mu dx^\mu$ is a connection 1-form.
$dA = F$ is the field strength 2-form (for the abelian case;
for non-abelian groups, $F = dA + A\wedge A$).
$\mathcal{F}_1[A] = \int|F|^2 = \SYM$.
The vacuum equation $d{*}F = 0$ is the source-free
Maxwell equation (abelian) or Yang--Mills equation
(non-abelian)~\known.

\smallskip
\noindent\textbf{0-form (Higgs):}
$H$ is a scalar in the fundamental representation of
$\mathrm{SU}(2)\times\mathrm{U}(1)$.
$\mathcal{F}_0[H] = \int|D_\mu H|^2$ with the gauge-covariant
derivative.
The vacuum equation $D^2H + V'(H) = 0$ includes the
potential term (spontaneous symmetry breaking), which
breaks the pure Fisher pattern.

\subsection{Unifying principle}

The key insight is \emph{not} that ``everything is $\Sigma$''---
it is that \emph{every sector of the Standard Model minimizes
Fisher information at its appropriate form degree.}
$\Sigma$ is the 0-form version (gravity); $A_\mu$ is the
1-form version (gauge).
The harmonic equation $d^*d\omega = 0$ is universal; what
differs is the form degree and the representation of the
gauge group.

This perspective clarifies the orthogonal
decomposition~\eqref{eq:decomposition}:
gravity and gauge forces are not different aspects of the
same scalar field $\Sigma$, but rather different entries in a
graded hierarchy of Fisher information minimization problems,
all generated by a single spectral action.


%=======================================================================
\section{Numerical verification}
\label{sec:numerical}
%=======================================================================

We verify the analytical results on the flat torus $T^d$
with equal periods $L_\mu = L$, where the Dirac spectrum
is known exactly~\known:
\begin{equation}\label{eq:torus_spectrum}
  \lambda_{\vec{n}}^2
  = \sum_{\mu=1}^{d}
  \left(\frac{2\pi n_\mu}{L}\right)^{\!2}\,,
  \qquad \vec{n}\in\mathbb{Z}^d\,.
\end{equation}

The CCSvS entropy~\cite{CCSvS2020} at conformal parameter~$Q$ is
\begin{equation}\label{eq:SvN_numerical}
  \SvN(Q)
  = \sum_{\vec{n}}\,
  s\!\left(\frac{\lambda_{\vec{n}}^2}{Q^2\Lambda^2}\right)\!,
\end{equation}
where
$s(x) = x\,\nF(x) + \ln(1+e^{-x})$ is the fermionic entropy
function with $\nF(x) = (e^x+1)^{-1}$.
The normalized Fisher metric is extracted via finite differencing:
\begin{equation}\label{eq:gQQ_numerical}
  g_{QQ}^{\mathrm{num}}
  = \frac{\SvN(Q+\delta)-2\SvN(Q)+\SvN(Q-\delta)}
  {\delta^2\,\SvN(1)}\,,
\end{equation}
with $\delta = 10^{-4}$ and evaluated at $Q=1$.

For the equal-period torus with $\Lambda L = 20$ (ensuring
$>10^4$ modes within the cutoff), we
obtain~\new:
\begin{equation}
  g_{QQ}^{\mathrm{num}}(1)\big|_{d=4}
  = 8.000 \pm 0.002
  = d \times (d-2)(d-3)/d = 4\times 2\,.
\end{equation}
Since all $d=4$ dimensions are rescaled simultaneously
under constant conformal rescaling,
$g_{QQ}/d = 2.000\pm 0.001$, confirming
Eq.~\eqref{eq:gQQ_EH} to better than $0.02\%$.

We also compute $g_{QQ}^{\mathrm{num}}$ on $T^d$ for
$d=2,3,4,5,6$.
Table~\ref{tab:numerical} confirms agreement with
$(d-2)(d-3)$ in every case.

\begin{table}[t]
\centering
\caption{Numerical verification of
$g_{QQ}^{\mathrm{EH}}(1) = (d-2)(d-3)$
on flat tori~$T^d$ with $\Lambda L = 20$.}
\label{tab:numerical}
\begin{ruledtabular}
\begin{tabular}{cccc}
$d$ & $(d\!-\!2)(d\!-\!3)$ & $g_{QQ}^{\mathrm{num}}/d$
  & Error \\
\colrule
2  &  0 & $0.000$ & --- \\
3  &  0 & $0.000$ & --- \\
4  &  2 & $2.000$ & $<0.02\%$ \\
5  &  6 & $6.00$  & $<0.1\%$ \\
6  & 12 & $12.0$  & $<0.2\%$ \\
\end{tabular}
\end{ruledtabular}
\end{table}


%=======================================================================
\section{What Doesn't Work}
\label{sec:limitations}
%=======================================================================

We state the limitations of the present construction explicitly.

\emph{(i) $F^2 \neq |\nabla\Sigma_{\mathrm{EM}}|^2$.}---The
gauge field strength $F_{\mu\nu}^a F^{a\,\mu\nu}$ is the
squared norm of a 2-form, \emph{not} of the gradient of a
scalar.
There exists no scalar field $\Sigma_{\mathrm{EM}}$ such that
$F_{\mu\nu}F^{\mu\nu} = |\nabla\Sigma_{\mathrm{EM}}|^2$
on a general 4-manifold.
The reason is topological: a 2-form $F$ can have non-trivial
cohomology class $[F]\in H^2(M;\mathbb{R})$, which obstructs
writing $F = d\Sigma_{\mathrm{EM}}$ (that would require $F$ to be
exact, hence cohomologically trivial).
The gauge sector is Fisher information of a \emph{1-form},
not a 0-form.
This is the fundamental reason why
$\Sgrav \oplus \Sgauge$ is an orthogonal
decomposition, not a unification into a single $\Sigma$.

\emph{(ii) The internal algebra $\AF$ is an input.}---The
finite spectral triple $F$ with algebra
$\AF = \mathbb{C}\oplus\mathbb{H}\oplus M_3(\mathbb{C})$
that generates the Standard Model gauge group
$\mathrm{SU}(3)\times\mathrm{SU}(2)\times\mathrm{U}(1)$
is not derived from the $\Sigma$ framework.
It is a separate input to the noncommutative
geometry~\cite{CCM2007,ConnesMarcolli2008}.
Whether $\AF$ can be constrained by information-theoretic
principles is an open question.

\emph{(iii) The cosmological constant term.}---The $a_0$
coefficient ($\propto\Vol(M)$) scales as $Q^d$ and has
Fisher metric $g_{QQ}^{(0)} = d(d-1)/Q^2 = 12/Q^2$ in
$d=4$.
This does not match $\Sigma = 2\ln Q$; it matches
$\Sigma = 12\ln Q$ (if such an identification were made).
The cosmological constant does not fit the Fisher framework
as a $\Sigma$-type quantity.

\emph{(iv) Fermion masses.}---The internal Dirac operator
$D_F$ encodes Yukawa couplings and fermion masses through the
finite spectral triple~\cite{CCM2007}.
This structure is entirely separate from the conformal
$\Sigma$ and is not addressed here.

\emph{(v) Constant conformal factor.}---Our scaling
analysis uses spatially uniform~$Q$.
For position-dependent $Q(x)$, the $a_4$ term develops
a non-trivial variation (the conformal anomaly), and the
simple formula $g_{QQ} = (d-2)(d-3)/Q^2$ acquires gradient
corrections~\cite{CC2006}.
The full result for general $Q(x)$ requires the
Chamseddine--Connes dilaton
construction~\cite{CC2006}, which we defer to future work.

\emph{(vi) Euclidean vs.\ Lorentzian signature.}---The
spectral action and CCSvS entropy are defined in Euclidean
signature.
Our $\Sigma = 2\ln Q$ is Lorentzian.
The Wick rotation is standard for static spacetimes but
not rigorously established for the full spectral
triple~\cite{vandenDungen2016,Besnard2018}.

\emph{(vii) The Khronon is not an inner fluctuation.}---In
the NCG framework, gauge fields arise as \emph{inner}
fluctuations of $D$~\cite{CC1997,vS2024}.
The Khronon field enters through \emph{outer}
automorphisms (conformal rescaling), not inner fluctuations.
The ghost condensation mechanism
$K(Q) = \mu^2(Q-1)^2$~\cite{BS2024} is additional
structure not part of the standard spectral triple.


%=======================================================================
\section{Discussion}
\label{sec:discussion}
%=======================================================================

\subsection{The ``why $d=4$'' result}

The strongest new contribution of this paper is the derivation
of $d=4$ from $(d-2)(d-3)=2$.
This is an \emph{algebraic} condition---not a dynamical
selection, not an anthropic observation, and not a
compactification from higher dimensions.
It states: the spacetime dimension is the unique value for
which the conformal Fisher information of the gravitational
action equals the entropy production of the quantum channel
that encodes temporal asymmetry.

The argument does not assume specific matter content (it holds
for any $\AF$), specific topology (it holds on any
closed manifold), or specific dynamics (it uses only
the heat kernel expansion, which is kinematic).
The only inputs are: (1)~the spectral action principle,
(2)~the identification $\Sigma = 2\ln Q$ from the Petz
recovery framework, and (3)~the requirement that these
two structures match at the level of Fisher information.

\subsection{Connection to Papers~I--VI}

Paper~I~\cite{PaperI} established $\Sigma = -\ln F \geq 0$
as a universal measure of temporal asymmetry.
Paper~II~\cite{PaperII} proved $\Sigma = 2\ln Q$ on static
backgrounds via the gravitational channel theorem.
Paper~III~\cite{PaperIII} extended this to FRW backgrounds
with $Q = 1+\delta$ and resolved the $c_s^2 = 0$ constraint.
Paper~VI~\cite{PaperVI} established the DBI-$\Sigma$
correspondence for electromagnetism.

The present paper identifies the origin of the coefficient
``2'' in $\Sigma = 2\ln Q$: it is $(d-2)(d-3)$ evaluated at
$d=4$.
It also explains why the gravitational and gauge sectors
decouple: they are Fisher information functionals at different
form degrees (0-form vs.\ 1-form), orthogonal under conformal
variation.

\subsection{The spectral action as entropy potential}

The CCSvS result~\cite{CCSvS2020} identifies the von~Neumann
entropy with the spectral action: $\SvN = \Tr\,h(\beta D)$.
Our result identifies $\Sigma$ with the Fisher information
(second variation) of this entropy.
The spectral action is therefore the ``entropy potential''
whose curvature in field space generates the temporal
asymmetry measure~$\Sigma$.

The Einstein--Hilbert term $a_2$ is the sector whose Fisher
curvature matches $\Sigma = 2\ln Q$; the matter sector $a_4$
is flat (conformally invariant) and does not contribute.
The hierarchy $a_0,\,a_2,\,a_4,\,a_6,\ldots$ is a hierarchy
of Fisher curvatures:
\begin{equation}
  g_{QQ}^{(k)}(1) = (d-2k)(d-2k-1)\big|_{d=4}
  = \{12,\, 2,\, 0,\, 6,\,\ldots\}\,,
\end{equation}
of which only $g_{QQ}^{(1)}=2$ matches $\Sigma$.

\subsection{The Dorau--Much connection}

Dorau and Much~\cite{DorauMuch2025} recently derived the
semiclassical Einstein equations from QRE on bifurcate
Killing horizons.
Their result corresponds to the \emph{first-order} variation
of $\Sigma$ (the entropic force $d\Sigma/dQ = 2/Q$), while
our Fisher metric is the \emph{second-order} variation.
The two are consistent: Einstein's equations arise from
$\delta\Sigma = 0$ at first order (equations of motion);
the Fisher metric characterizes the information geometry
of the solution space at second order (fluctuations around
equilibrium).

\subsection{Towards the Standard Model from $\Sigma$}

The conformal invariance of $a_4$ in $d=4$ means that gauge
fields and the Higgs potential do not affect $\Sigma$ at the
level of constant conformal variation.
The gauge sector enters through \emph{position-dependent}
variations (inner fluctuations of $D$), which generate the
Yang--Mills and Higgs equations as Euler--Lagrange conditions
of the spectral action~\cite{CC1997,vS2024}.
A full unification requires extending the present analysis
from constant to general conformal factors, incorporating
both inner and outer fluctuations.
This is the subject of ongoing work (Papers~7--8).

\subsection{Open questions}

Several questions remain:
\begin{enumerate}
\item Can the internal algebra $\AF$ be derived from
  $\Sigma$ or from information-theoretic principles?
\item Does the position-dependent extension $Q(x)$ modify
  the dimension selection, or is $(d-2)(d-3)=2$ robust?
\item Can fermion masses (encoded in $D_F$) be connected
  to the Fisher framework?
\item Is there a Lorentzian spectral triple that makes the
  Wick rotation unnecessary?
\end{enumerate}
These are non-trivial open problems.
The present paper establishes the foundation: the spectral
action and the $\Sigma$ framework are connected through
Fisher information, and this connection selects $d=4$.


%=======================================================================
\begin{acknowledgments}
The author thanks Mark~M.~Wilde for correspondence
on the Petz recovery map and acknowledges the foundational
work of Alain Connes, Ali~H.~Chamseddine, and
Walter~D.~van~Suijlekom on the spectral action and its
entropic interpretation.
\end{acknowledgments}
%=======================================================================


\begin{thebibliography}{99}

\bibitem{PaperI}
S.-K.~Huang,
``Temporal asymmetry from the Petz recovery map:
A unifying measure for the arrow of time,''
arXiv:2502.xxxxx (2026).

\bibitem{PaperII}
S.-K.~Huang,
``The gravitational refractive index as Petz recovery
fidelity: Temporal asymmetry from spacetime geometry,''
in preparation (2026).

\bibitem{PaperIII}
S.-K.~Huang,
``Dark matter as Khronon condensate: Weak-field
phenomenology of the $\Sigma$ framework,''
in preparation (2026).

\bibitem{PaperVI}
S.-K.~Huang,
``Electromagnetic force from temporal asymmetry: U(1)
gauge structure, DBI-$\Sigma$ correspondence, and
modified Maxwell equations,''
in preparation (2026).

\bibitem{Ehrenfest1917}
P.~Ehrenfest,
``In what way does it become manifest in the fundamental
laws of physics that space has three dimensions?''
Proc.\ Amsterdam Acad.\ \textbf{20}, 200 (1917).

\bibitem{Polchinski1998}
J.~Polchinski,
\emph{String Theory}, Vols.\ 1--2
(Cambridge University Press, Cambridge, 1998).

\bibitem{CC1997}
A.~H.~Chamseddine and A.~Connes,
``The spectral action principle,''
Commun.\ Math.\ Phys.\ \textbf{186}, 731 (1997);
hep-th/9606001.

\bibitem{CCM2007}
A.~H.~Chamseddine, A.~Connes, and M.~Marcolli,
``Gravity and the standard model with neutrino mixing,''
Adv.\ Theor.\ Math.\ Phys.\ \textbf{11}, 991 (2007);
hep-th/0610241.

\bibitem{ConnesMarcolli2008}
A.~Connes and M.~Marcolli,
\emph{Noncommutative Geometry, Quantum Fields and Motives}
(American Mathematical Society, Providence, 2008).

\bibitem{Connes1996}
A.~Connes,
``Gravity coupled with matter and the foundation of
non-commutative geometry,''
Commun.\ Math.\ Phys.\ \textbf{182}, 155 (1996);
hep-th/9603053.

\bibitem{vS2024}
W.~D.~van~Suijlekom,
\emph{Noncommutative Geometry and Particle Physics},
2nd ed.\ (Springer, Cham, 2024).

\bibitem{CCSvS2020}
A.~H.~Chamseddine, A.~Connes, and
W.~D.~van~Suijlekom,
``Entropy and the spectral action,''
Commun.\ Math.\ Phys.\ \textbf{373}, 457 (2020);
arXiv:1809.02944.

\bibitem{Vassilevich2003}
D.~V.~Vassilevich,
``Heat kernel expansion: User's manual,''
Phys.\ Rep.\ \textbf{388}, 279 (2003);
hep-th/0306138.

\bibitem{Gilkey1995}
P.~B.~Gilkey,
\emph{Invariance Theory, the Heat Equation, and
the Atiyah--Singer Index Theorem},
2nd ed.\ (CRC Press, Boca Raton, 1995).

\bibitem{BGV2004}
N.~Berline, E.~Getzler, and M.~Vergne,
\emph{Heat Kernels and Dirac Operators},
corrected reprint (Springer, Berlin, 2004).

\bibitem{Wald1984}
R.~M.~Wald,
\emph{General Relativity}
(University of Chicago Press, Chicago, 1984).

\bibitem{CC2006}
A.~H.~Chamseddine and A.~Connes,
``Scale invariance in the spectral action,''
J.\ Math.\ Phys.\ \textbf{47}, 063504 (2006).

\bibitem{DorauMuch2025}
A.~Dorau and A.~Much,
``From quantum relative entropy to the semiclassical
Einstein equations,''
Phys.\ Rev.\ Lett.\ (2025);
arXiv:2510.24491.

\bibitem{BS2024}
L.~Blanchet and C.~Skordis,
``Dark matter as a Khronon condensate,''
arXiv:2404.06584 (2024).

\bibitem{Frieden2004}
B.~R.~Frieden,
\emph{Science from Fisher Information},
2nd ed.\ (Cambridge University Press, Cambridge, 2004).

\bibitem{BaezHuerta2010}
J.~C.~Baez and J.~Huerta,
``Division algebras and supersymmetry~I,''
in \emph{Superstrings, Geometry, Topology,
and $C^*$-Algebras}, Proc.\ Symp.\ Pure Math.\
\textbf{81}, 65 (2010);
arXiv:0909.0551.

\bibitem{Jacobson1995}
T.~Jacobson,
``Thermodynamics of spacetime: The Einstein equation
of state,''
Phys.\ Rev.\ Lett.\ \textbf{75}, 1260 (1995);
gr-qc/9504004.

\bibitem{vandenDungen2016}
K.~van~den~Dungen,
``Krein spectral triples and the fermionic action,''
Math.\ Phys.\ Anal.\ Geom.\ \textbf{19}, 4 (2016).

\bibitem{Besnard2018}
F.~Besnard,
``On the definition of spacetimes in noncommutative
geometry: Part~I,''
J.\ Geom.\ Phys.\ \textbf{123}, 292 (2018);
arXiv:1611.07830.

\bibitem{Umegaki1962}
H.~Umegaki,
``Conditional expectation in an operator algebra.
IV. Entropy and information,''
K\=odai Math.\ Sem.\ Rep.\ \textbf{14}, 59 (1962).

\end{thebibliography}

\end{document}
